\documentclass[12pt,a4paper]{scrreprt}
\usepackage[utf8]{inputenc}
\usepackage[T1]{fontenc}
\usepackage{lmodern}
\usepackage{graphicx}
\usepackage{amsmath, amssymb}
\usepackage{natbib}
\usepackage{hyperref}
\bibliographystyle{apalike}
\begin{document}

\section*{Categorical Time-Series Analysis}

\subsection*{General Notation}

Let $(X_t)$ be a \textit{categorical}  time series process, where each $X_t$ takes values in
\[
S := \{s_0, \dots, s_m\}, \qquad m \in \mathbb{N}.
\]

$(X_t)$ may be \textit{nominal}, in which case the order of magnitude of the elements in $S$ is arbitrary and solely usesd for notational consistency, or $ordinal$, in which case the order of the elements in $S$ reflects the inherent order of the elements.\\  

Furthermore, $(X_t)$ is assumed to be stationary. A distinction is made between two forms of stationarity. A process is said to be weakly stationary if:
    \[
\begin{aligned}
E(X_t) &= \mu_X \quad &&\forall t, \\
E(X_t^2) &< \infty \quad &&\forall t, \\
\mathrm{Cov}(X_t, X_{t+h}) &= \gamma(h) \quad &&\forall t,h. 
\end{aligned}
\]

In contrast, strong stationarity requires that all finite-dimensional distributions are invariant under time shifts, i.e.,
\[F_{X_t, \ldots, X_{t + j}} = F_{X_{t+h}, \ldots, X_{t+h + j}} \quad \forall t, j, h\]
\subsection*{Marginal Distributions}

Regardless of whether we consider a simple categorical process with few categories or attempt to describe a complex, temporally evolving process using moments and measures of serial dependence, the analysis of categorical processes is always built on the examination of the marginal probabilities. However, because ordinal and nominal data encode different types of information, their marginals are typically represented in different ways:

\begin{itemize}
    \item \textbf{Ordinal data:}  
    Analytical tools for ordinal processes must preserve the natural ordering of the elements in $S$. Measures that ignore this ordering—such as the mode—would discard essential information about the relative positions of the categories. For this reason, methods for ordinal data rely on the marginal cumulative distribution function (CDF), which accumulates probabilities in accordance with the natural ordering of the categories, thereby preserving the ordinal structure, and is therefore the standard representation for ordinal data:
    \[
        \mathbf{f} = (f_0, \dots, f_{m-1})^\top,
        \qquad 
        f_i := P(X_t \le s_i).
    \]

    For ordinal data, it is sufficient to consider the cumulative probabilities of the categories $0, \ldots , m-1$ since the cumulative probability of the last category $m$ is always equal to 1. \\
    
    Although one could equivalently estimate the PMF and derive the CDF from it, representing ordinal data through the CDF is often more convenient and aligns with widely used modelling frameworks.

    \item \textbf{Nominal data:}  
    Since the categories of nominal variables do not possess a meaningful order, statistical procedures cannot make use of cumulative structure.  
    Therefore the marginal distribution is best represented by the probability mass function (PMF):
    \[
        \mathbf{p} = (p_0, \dots, p_{m})^\top,
        \qquad 
        p_i := P(X_t = s_i).
    \]
\end{itemize}


\subsection*{Bivariate Probabilities at Lag $h$}

Temporal dependence can be described by bivariate (lag-$h$) probabilities:

\begin{itemize}
    \item \textbf{Ordinal series (cumulative):}
    \[
    f_{ij}(h) := P(X_t \le s_i,\, X_{t+h} \le s_j), \qquad i,j < m.
    \]

    \item \textbf{Nominal series (joint):}
    \[
    p_{ij}(h) := P(X_t = s_i,\, X_{t+h} = s_j), \qquad i,j \le m.
    \]
\end{itemize}

These quantities serve as input for dependence measures such as association coefficients or transition structures.
\subsection*{Location}

The choice of a location measure for a categorical process depends on whether the process is ordinal or nominal. 
For an ordinal process, the natural ordering of the categories must be respected, and the median serves as an appropriate measure of central tendency:
\[
\operatorname{Med}(X_t)
=
\min \bigl\{\, s_i \in S : F(s_i) \ge \tfrac{1}{2} \,\bigr\}.
\]

For a nominal process, any measure that relies on ordering is inappropriate. 
Instead, a location measure must depend only on the category probabilities. 
The mode, defined as the most probable category,
\[
\operatorname{Mod}(X_t) = \arg\max_{s_i \in S} P(X_t = s_i),
\]
is therefore the natural choice.
\subsection*{Measures of Dispersion}

The concept of dispersion differs between ordinal and nominal processes. Therefore, two separate measures are used: the \emph{Index of Ordinal Variation} (IOV) for ordinal data and the \emph{Index of Qualitative Variation} (IQV) for nominal data. Both indices are scaled to the interval $[0,1]$, where a value of $0$ represents minimal dispersion and a value of $1$ represents maximal dispersion.\\


\[
\mathrm{IOV}
= \frac{4}{m}\sum_{i=0}^{m-1} f_i (1 - f_i) \qquad \mathrm{IQV} = \frac{m+1}{m}\bigg(1 - \sum_{i=0}^{m} p_i^2\bigg) \]


\subsubsection*{Minimal Dispersion}

Both measures attain the value $0$ if and only if the distribution is 
degenerate, that is, if all probability mass is concentrated in a single 
category.

\paragraph{Ordinal case.}
A one–point distribution assigns all probability mass to a single category 
$s_j \in S$.  
For ordinal data, this corresponds to a CDF vector $\mathbf{f}$ of the form
\[
\mathbf{f} \in 
\left\{
\begin{pmatrix} 1 \\ 1 \\ \vdots \\ 1 \end{pmatrix},
\begin{pmatrix} 0 \\ 1 \\ \vdots \\ 1 \end{pmatrix},
\ldots,
\begin{pmatrix} 0 \\ 0 \\ \vdots \\ 1 \end{pmatrix}
\right\}
:= \{\boldsymbol{c}_0,\ldots,\boldsymbol{c}_m\} \subset [0,1]^{m+1}.
\]
In each of these cases, every component satisfies $f_i \in \{0,1\}$, and thus
\[
f_i(1 - f_i) = 0 \qquad \forall i.
\]
Therefore,
\[
\mathrm{IOV}
= \frac{4}{m}\sum_{i=0}^{m-1} f_i (1 - f_i)
= 0.
\]
Conversely, if $\mathrm{IOV}=0$, then
\[
\sum_{i=0}^{m-1} f_i(1 - f_i) = 0,
\]
and since each term is nonnegative,
\[
f_i(1 - f_i) \overset{!}{=} 0 \quad \forall i.
\]
This is only possible if $f_i \in \{0,1\}$ for all $i$, which implies
$\mathbf{f} \in \{\boldsymbol{e}_0,\ldots,\boldsymbol{e}_m\}$.  
Hence, $\mathrm{IOV}=0$ holds \emph{exactly} for one–point distributions.
\paragraph{Nominal case.}
A degenerate (one--point) distribution assigns all probability mass to a single 
category $s_j \in S$, which corresponds to a PMF vector of the form
\[
\mathbf{p} \in 
\left\{
\begin{pmatrix} 1 \\ 0 \\ \vdots \\ 0 \end{pmatrix},
\begin{pmatrix} 0 \\ 1 \\ \vdots \\ 0 \end{pmatrix},
\ldots,
\begin{pmatrix} 0 \\ 0 \\ \vdots \\ 1 \end{pmatrix}
\right\}
:= \{\boldsymbol{e}_0,\ldots,\boldsymbol{e}_m\} \subset [0,1]^{m+1}..
\]
In each of these cases, exactly one component equals $1$ and all others equal $0$, so
\[
\sum_{i=0}^m p_i^2 = 1.
\]
Thus,
\[
\mathrm{IQV} 
= \frac{m+1}{m}\bigl(1 - \sum_{i=0}^m p_i^2\bigr) 
= \frac{m+1}{m}(1 - 1) = 0.
\]
Conversely, if $\mathrm{IQV} = 0$, then
\[
\frac{m+1}{m}\Bigl(1 - \sum_{i=0}^m p_i^2\Bigr) = 0
\quad\Longrightarrow\quad
\sum_{i=0}^m p_i^2 = 1.
\]
As all probabilities satisfy $0 \leq p_i \leq 1$, it follows that \[
p_i^2 \leq p_i \quad \forall p_i, \text{ and } p_i^2 < p_i \quad \forall p_i \notin \{0,1\}.
\]
Since $\sum_{i=1}^{m}p_i = 1$, it follows that the only way $\sum_i p_i^2 = 1$ can hold is if exactly one $p_i$ equals 1 and all others are 0. In other words, the distribution is degenerate and $\boldsymbol{p} \in \{\boldsymbol{e}_0,\ldots, \boldsymbol{e}_{m}\}$.

 


\subsubsection*{Maximal Dispersion}

\paragraph{Ordinal case.}
Dispersion of an ordinal variable is maximal when half of the 
probability mass lies in the lowest category and half in the highest 
category; with all intermediate categories having probability zero.  
This results in a CDF vector,

\[
\mathbf{f} = \left(\tfrac{1}{2}, \tfrac{1}{2}, \ldots, \tfrac{1}{2}\right)^\top.
\]
Then
\[
f_i(1 - f_i) = \frac{1}{2}\left(1 - \frac{1}{2}\right) = \frac{1}{4} \qquad \forall i,
\]
and thus
\[
\mathrm{IOV}
= \frac{4}{m} \cdot \sum_{i=0}^{m-1} \frac{1}{4}
= \frac{4}{m} \cdot \frac{m}{4}
= 1.
\]
To show that $\mathrm{IOV} = 1 \Rightarrow \boldsymbol{f} = (\tfrac{1}{2}, \ldots, \tfrac{1}{2})^T$, recall that the IOV only depends on $\boldsymbol{f}$ through the sum
\[
\sum_{i=0}^{m-1} f_i (1 - f_i).
\]
Taking the derivative of $g(f_i) := f_i(1-f_i)$  and setting it to zero yields:
\begin{align*}
\frac{\partial}{\partial f_i} g(f_i) &= 1 - 2 f_i, \\
1 - 2 f_i &\overset{!}{=} 0, \\
f_i &= \frac{1}{2}.
\end{align*}

Since $g(f_i=0) = g(f_i = 1)=0$, the fucntion is concave and this critical point is indeed the global maximum.  
Since its been already shown that $IOV(\boldsymbol{f} = (\frac{1}{2}, \ldots, \frac{1}{2})^T) = 1$, it directly follows that
\[
IOV(\boldsymbol{f}) = 1 \Leftrightarrow \boldsymbol{f} = (\frac{1}{2}, \ldots, \frac{1}{2})^T
\]


\paragraph{Nominal case.}
Maximal dispersion is obtained when probability mass is distributed 
uniformly across all $m+1$ categories:
\[
p_i = \frac{1}{m+1}, \qquad i = 0,\ldots,m.
\]
Then
\[
\sum_{i=0}^m p_i^2 = (m+1)\frac{1}{(m+1)^2}
= \frac{1}{m+1},
\]
so that
\[
\mathrm{IQV}
= \frac{m+1}{m}\left(1 - \frac{1}{m+1}\right)
= \frac{m+1}{m} \cdot \frac{m}{m+1}
= 1.
\]
And if $\mathrm{IQV} = 1$:
\[\frac{m+1}{m} (1- \sum_{i=0}^{m} p_i^2) = 1 \Rightarrow (1-\sum_{i=1}^{m}p_i^2) = \frac{m}{m+1} \Rightarrow \sum_{i=1}^{m} p_i^2 = \frac{1}{m+1}\]
Given $\sum_{i=0}^{m}p_i = 1$ it follows that
\[\arg \underset{\boldsymbol{p}}{\min} \sum_{i=0}^{m} p_i^2 = (\frac{1}{m+1}, \ldots, \frac{1}{m+1})^T\] 
since it generally holds that given some vector $X$ and its components sum $\sum_{i=1}^{m}x_i = c$, with c being some constant, that 
\[\arg \underset{\boldsymbol{X}}{\min} \sum_{i=1}^{m}x_i = (\frac{1}{m}, \ldots, \frac{1}{m}) \]
\section*{Samples Measures}

For the estimation of the (cumulative) probabailities based on samples we use the (cumulative) relative frequencies, that can be be best expressed if switching from $(X_t)$ to its binarization ($(\boldsymbol{Z_t})$ for ordinal variables, $(\boldsymbol{Y}_t)$ for nominal):
\[\boldsymbol{Z}_t = (Z_{t, 0}, \ldots, Z_{t, m-1})^T,  Z_{t, i} = I_{\{X_t \leq s_i\}} \qquad \boldsymbol{Y}_t = (Y_{t, 0}, \ldots, Y_{t, m})^T, Y_{t, i} = I_{\{X_t = s_i\}}  \]

Now we can estimate $\boldsymbol{p}/ \boldsymbol{f}$ by taking the mean about the binarization:
\[\boldsymbol{\hat{f}} = \frac{1}{n}\sum_{t=1}^{n} \boldsymbol{Z}_t, \quad \boldsymbol{\hat{p}} = \frac{1}{n}\sum_{t=1}^{n} \boldsymbol{Y}_t \]

To get the asymptotics of the different sample measures, we first need a CLT on $ \boldsymbol{\hat{f}}/ \boldsymbol{\hat{p}}$ \\
,0
As our data are not (necesarilly) indepedent, the CLT only holds for a stationary process that fullfills a strong mixing condition (e.g. alpha mixing; Bradley 2005). The stationarity of $\boldsymbol{Z_t}/ \boldsymbol{Y_t}$ follows from the stationarity of $X_t$, as the transforming function is deterministic. \\

A sequence is said to be $\alpha$-mixing, if: 
\begin{align*}
    \text{For any event } &A \in \sigma(X_s : s \leq t) \text{ and event } B \in \sigma(X_s : s \geq t+h ) \text{ the }\alpha\text{-coefficient} \\
    & \alpha(h)=  \sup |P(A\cup B) - P(A)P(B)| \text{ goes to } 0 \text{ for } h \rightarrow \infty
\end{align*}

If $\alpha$-mixing is satisfied it holds that: 
 
\[\sqrt{n} (\overbrace{\bar{\boldsymbol{Z}}_t}^{= \boldsymbol{\hat{f}}}- \boldsymbol{f} )\overset{d}{\rightarrow} \mathcal{N}(0, \boldsymbol{\Sigma}), \quad\text{ with } \boldsymbol{\Sigma} = (\sigma_{i,j})_{i,j = 0, ~ \ldots, m-1} \]
\[\text{where }\sigma_{i,j} = Cov(Z_{0,i},Z_{0,j}) + \sum_{h=1}^{\infty} \left(Cov(Z_{0, i}, Z_{h, j}) + Cov(Z_{h, i}, Z_{0, j})\right) \]
\[\text{and } Cov(Z_{t, i}, Z_{s, j}) = \begin{cases}
    f_{ij}(t-s)-f_i f_j & \text{if } t \geq s \\
    f_{ji}(s-t)-f_i f_j & \text{if } s \geq t 
\end{cases}\]

To now derive the asymptotics of the dispersion measures (scalar functions) we can apply the delta method (i.e. Theorem on 5.25 on Asymptotic Distribution of $g(\boldsymbol{X}_n)$ in the formulary) according to which the following holds:
 
\[n^{\frac{1}{2}}(IOV(\boldsymbol{\hat{f}})- IOV(\boldsymbol{f})) \overset{d}{\rightarrow} \mathcal{N}(0, \boldsymbol{G\Sigma G})\]
\[\text{Where } G_{1 \times m} = \left[\frac{\partial IOV(\boldsymbol{f})}{\partial \hat{f}_0}, \ldots, \frac{\partial IOV(\boldsymbol{f})}{\partial \hat{f}_{m-1}}\right]' \]

The derivations for other dispersion metrics are analogous.

\section*{Modelling Missingness}

To model the missingness amplitude modulation is used. The concept stems from physics and allows to model a stochastic process that runs concurrent to another one an modulates its amplitude. In the case of missingness the main process $(X_t)_{t}$ is either there or not there. Therefor we can define our apmlitude modulating process as 
\[O_t \in \{0,1\} \quad \forall t\] 
where $O_t = 1$ if $(X_t)$ is observed and 0 if it isn't

\end{document}
